\documentclass[10pt]{article}

\usepackage[utf8]{inputenc}

\usepackage[T1]{fontenc}

\usepackage{amsmath}

\usepackage{amsfonts}

\usepackage{amssymb}

\usepackage[version=4]{mhchem}

\usepackage{stmaryrd}

\usepackage{bbold}

\usepackage{graphicx}

\usepackage[export]{adjustbox}

\graphicspath{ {./images/} }



\begin{document}

Для молекул имеются специфич. О. П. для электронных, колебательных и вращательных молекулярных спектров, определяемые симметрией равновесных конфигураций молекул, а для кристаллов - О.п. для их электронных и колебательных спектров, определяемые симметрией кристаллич. решетки.



В физике элементарных частиц, кроме общих законов сохранения энергии, импульса, момента импульса, имеются дополнительные законы сохранения, связанные с симметриями фундаментальных взаимодействий частиц - сильного, электромагнитного и слабого. Процессы превращений элементарных частиц подчиняются строгим законам сохранения электрич. заряда $Q$, барионного заряда $B$ и лептонного заряда $L$, к-рым соответствуют строгие О. п.: $\Delta Q=$ $=\Delta B=\Delta L=0$ (это означает, что при достигнутой точности измерений нарушения этих О. П. не обнаружены).



М. А. Ельяшевич.



OTКРЫТАЯ ВСЕЛЕННАЯ - см. КосМологИческИе Модели. ОТКРЫТАЯ КВАНТОВАЯ СИСТЕМА - см. Квантовая открытая система.



ОTKРЫTOE МНОЖЕСТВО - см. Топологическое пространство.



ОТНОСИТЕЛЬНАЯ ТУРБУЛЕНТНАЯ ДИФФУЗИЯ см. Ричардсона закон четырех третей.



ОтНОСИТЕЛЬНОЕ РАВНОВЕСИЕ - состояние гамильтоновой системы, которому соответствует фазовая кривая, проектирующаяся в положение равновесия приведенной системы.



Пусть гамильтонова система с гамильтонианом $H$ обладает группой с симметрией $G$. Симплектическая редукция позволяет рассмотреть приведенную систему на приведенном фазовом пространстве $F_{p}$. Фазовые кривые исходной системы $\boldsymbol{M}$, проектирующиеся в положение равновесия приведенной системы, отвечают положениям О.p.



Фазовая кривая соответствует О. р. тогда и только тогда, когда она является орбитой однопараметрич. подгруппы группы $\boldsymbol{G}$ в исходном пространстве. Другое описание О. p.критич. точки отображения момента и энергии



\[

P \times H: M \rightarrow \mathbb{S}^{*} \times \mathbb{R}

\]



на регулярном множестве уровня момента. $\quad$ М. Э. Казарян.



\includegraphics[max width=\textwidth, center]{2023_07_08_1faf033065061aa7d0b4g-1}

твердого тела - вид движения твердого тела, подчиненного нестационарной связи, при котором тело не изменяет своего положения относительно этой связи. Обычно нестационарная связь реализуется путем принудительного перманентного вращения нек-рого второго тела. Тогда О.р. проявляется в том, что оба тела совершают перманентное вращение как одно целое. Напр., горизонтальная ось качания физич. маятника совершает принудительное вращение вокруг вертикали с постоянной угловой скоростью. В этом случае сушествует несколько положений О.р. маятника. Нек-рые из них могут изменяться вместе с изменением величины угловой скорости вращения оси.



См. также Равновесие механической системы.



В. А. Самсонов.



ОтНОСИтЕЛЬНОСтИ ПРИНЦИП, принци п относ и т е л ьн о с т и Эйн ш те йна - утверждение, состоящее в том, что все физические явления (механические, оптические, электромагнитные и любые другие) при всех одинаковых начальных условиях протекают одинаково во всех инерииальных системах отсчета. Этот постулат был, по-видимому, впервые высказан А. Пуанкаре (Н. Poincaré, 1895). Вместе с постулатом о независимости скорости света от движения источника О. п. был положен А. Эйнштейном (A. Einstein) в основу построения относительности теории, приведшей к глубокому пересмотру понятий о пространстве и времени. О.П. содержит Галилея принцип относительности как предельный случай при малых по сравнению со скоростью света скоростях тел. И. Ю. Кобзарев. ОТНОСИТЕЛЬНОСТИ ТЕОРИЯ - физическая теория, рассматривающая пространственно-временные закономерности, справедливые для любых физических процессов. Универсальность пространственно-временных свойств, рассматриваемых О. т., позволяет говорить о них просто как о свойствах пространства-времени. Наиболее общая теория пространства-времени называется о бщ е й те о р и е й относительности, или теорией тяготения, так как согласно этой теории свойства пространства-времени в данной области определяются действующими в ней полями тяготения. В частной теори и относительности, основы к-рой были даны А. Эйнштейном (А. Einstein, 1905), изучаются свойства пространства-времени, справедливые с той точностью, с какой можно пренебрегать действием тяготения. Таким образом, логически частная О.т. есть частный случай общей О.т.; построение общей О.т. было завершено А. Эйнштейном позже (в 1915), после чего и появился термин «частная О.Т.». В русской литературе последняя называется также с пеци альн ой те о р и ей относительности или просто О.т.



Лит.: [1] Математическая энциклопедия, т. 4, М., 1984, с. 145-51; [2] Физическая энциклопедия, т. 3, М., 1992, с. 493-502; см. также лит. к статьям в [1] и [2].



ОТНОСИТЕЛЬНОСТЬ ОДНОВРЕМЕННОСТИ - См. Peлятивистские эффекты.



ОТНОСИТЕЛЬНЫЙ ИНВАРИАНТ - см. Вейеритрасса эллиптические функции, Модулярная функиия.



ОТНОСИТЕЛЬНЫЙ ИНТЕГРАЛЬНЫЙ ИНВАРИАНТ дифференциальная $k$-форма $\omega$ на гладком многообразии $M$, интегралы от которой по замкнутым $k$-мерным цепям сохраняются заданным диффеоморфизмом



\[

g: M \rightarrow M, \int_{g c} \omega=\int_{c} \omega .

\]



Если $\omega-$ O. и. и., то $d \omega-$ абсолютный интегральный инвариант. Обратное верно, только если любая замкнутая $k$-мерная цепь в $M$ является границей. См. Интегральный инвариант.



П р и м е p. Канонич. отображение $g: \mathbb{R}^{2 n} \rightarrow \mathbb{R}^{2 n}$ имеет относительным инвариантом 1-форму $p d g=\sum p_{i} d q_{i}$.



М. Э. Казарян.



ОТОБРАХЕНИЕ ПЕРИОДОВ ДИфференциальНОЙ форМЫ $\omega$ на пространстве гладкого расслоения - многозначное отображение базы в когомологии отмеченного слоя. Точке $\lambda$ базы ставится в соответствие класс когомологий формы $\omega$, перенесенный Гаусса - Манина связностью в отмеченный слой. Если $\delta_{1}(\lambda), \ldots, \delta_{\mu}(\lambda)-$ ковариантно постоянный базис гомологий, то в этих координатах О. п. имеет вид:



\[

\lambda \rightarrow\left(\int_{\delta_{1}(\lambda)} \omega, \ldots, \int_{\delta_{\mu}(\lambda)} \omega\right)

\]



С. В. Чмутов.



ОТОБРАЖЕНИЙ МЕТОД - см. Изображений метод. ОТРАХЕНИЙ ТРУППА - см. Вращений группа.



ОТРИЦАТЕЛЬНАЯ ВЯЗКОСТЬ - Явление, наблюдающееся в турбулентно движущихся жидкостях (газах, плазме) и заключающееся в нелинейном переносе кинетической энергии мелкомасштабных (турбулентных) движений среды к крупномасштабным (осредненным). При предположении о пропорциональности тензора напряжений Рейнольдса тензору скоростей деформации осредненного поля скорости среды, коэффициенты пропорциональности (коэффициенты турбулентной вязкости) в случае явления О. в. оказываются отрицательными. Явление О. в. обнаружено экспериментально в атмосфере Земли, океане, фотосфере Солнца.



Лит.: [1] С тар р В., Физика явлений с отрицательной вязкостью, пер. с англ., М., 1971.



Р. В. Озмидов.



ОТТАЛКИВАЮЩЕЕ МНОГООБРАЗИЕ - сМ. Центральное многообразие.



OЧАРОВАНИЕ квантового ч исла-см. Аромат. ОШИБКА - см. Анализ данных.



ОШИБОК ИНТЕГPAЛ - см. Интегральные функции.





\end{document}